\documentclass[UTF8,a4paper,12pt]{ctexart}

% 2026 华为杯全国研究生数学建模竞赛 LaTeX 模板（简化版）
% 编译方式：XeLaTeX

\usepackage{geometry}
\geometry{left=2.5cm,right=2.5cm,top=2.5cm,bottom=2.5cm}

\usepackage{amsmath,amssymb,amsfonts}
\usepackage{graphicx}
\usepackage{booktabs}
\usepackage{algorithm}
\usepackage{algorithmic}
\usepackage{hyperref}
\usepackage{enumitem}

\title{\textbf{论文题目}}
\author{队伍编号：XXXXXX}
\date{}

\begin{document}

\maketitle

\begin{abstract}
本文针对XXXX问题，建立XXXX模型，提出XXXX方法。
关键词：数学建模；优化模型；机器学习；预测
\end{abstract}

\tableofcontents
\newpage

\section{问题重述}
介绍赛题背景和需要解决的问题。

\section{问题分析}

\subsection{总体思路}
给出整体建模流程。

\begin{figure}[htbp]
\centering
\IfFileExists{figures/framework.png}{%
  \includegraphics[width=0.75\linewidth]{figures/framework.png}%
}{%
  \fbox{\parbox[c][3cm][c]{0.7\linewidth}{\centering
    在此插入整体技术路线图\\[0.5em]
    \texttt{figures/framework.png}}}%
}
\caption{整体技术路线}
\end{figure}

\section{模型假设}
\begin{enumerate}
\item 假设1；
\item 假设2；
\item 假设3。
\end{enumerate}

\section{符号说明}

\begin{table}[htbp]
\centering
\begin{tabular}{ccc}
\toprule
符号 & 含义 & 单位\\
\midrule
$x$ & 示例变量 & -\\
$y$ & 输出变量 & -\\
\bottomrule
\end{tabular}
\caption{主要符号说明}
\end{table}

\section{模型建立}

\subsection{模型一}
数学公式：

\[
\min f(x)=\sum_{i=1}^{n}(y_i-\hat y_i)^2
\]

\subsection{算法设计}

\begin{algorithm}
\caption{算法流程}
\begin{algorithmic}
\STATE 初始化参数
\STATE 输入数据
\STATE 迭代优化
\STATE 输出结果
\end{algorithmic}
\end{algorithm}

\section{模型求解}

介绍求解过程、代码实现和实验结果。

\section{结果分析}

\begin{table}[htbp]
\centering
\begin{tabular}{ccc}
\toprule
方法 & 指标1 & 指标2\\
\midrule
Model A & 0.95 & 10\\
Model B & 0.97 & 8\\
\bottomrule
\end{tabular}
\caption{模型性能比较}
\end{table}

\section{模型评价}

分析优点、不足和改进方向。

\section{参考文献}

\bibliographystyle{plain}
% \bibliography{reference}

\appendix
\section{附录}

附录代码、补充证明等。

\end{document}
